\section{引言与相关工作}\label{sec:intro}

室内传感与智能家居网络常在 \SI{2.4}{GHz} 频段并存部署低功耗蓝牙（\ble{}）~\cite{bluetooth_core51} 与 \stdfour{}~\cite{ieee802154_2020} 节点。在没有多电台网关的情况下，单一发送机向两类协议设备下发控制命令通常只能按协议顺序分别发送突发~\cite{brunner2021leveraging}，而超出 \SI{31}{B} 传统 \ble{} 广播净荷上限的命令还必须跨多帧分片。无连接 \ble{}~5.1 扩展广播将至多 \SI{255}{B} 卸载至次信道 \texttt{AUX\_ADV\_IND} 协议数据单元（PDU）~\cite{bluetooth_core51}，从而使单个次信道突发能够同时承载独立命令与共享时间元数据。

既有跨技术通信（CTC）研究已覆盖物理层帧仿真（BlueBee~\cite{jiang2017bluebee} 将 \stdfour{} 码片映射至 \SI{1}{Mbps} \ble{} 高斯频移键控（GFSK）比特，带来 $4\times$ 字节膨胀）、包级异构广播调度（Brunner 等人~\cite{brunner2021leveraging}）、宽带 IEEE~802.11ax 正交频分多址（OFDMA）并发传输（WiCast~\cite{xia2025wicast}），以及基于外部参考的物联网（IoT）设备双向时钟同步（X-Sync~\cite{hofmann2021xsync}、TimeBee~\cite{gao2023timebee}）。然而，针对不具备宽带 OFDMA 能力的单个 \SI{2}{MHz} 带宽 \ble{}~5 广播射频，既有工作尚未定量刻画：当单次任务要求\emph{双侧接收链均校验通过}各自独立命令时，共包封装开销、等预算重发行为以及传统广播分片暴露风险如何制约整体权衡。在不把双协议广播宣称为首次提出、不提出跨设备时钟同步算法、也不将本地预测误差与 X-Sync~\cite{hofmann2021xsync} 排名的前提下，本文研究固定次信道单个 \texttt{AUX\_ADV\_IND} 共包何时能够在提交波形成本与计划尝试内联合解析率上取得收益。

为在匹配的逻辑任务语义下回答该问题，本文设计了统一紧凑双协议封装（\proto{}）及槽隔离传统分片基线（\fragproto{}），并在单台软件定义无线电（SDR）发送机与两套独立 SDR 接收链上实现评估。本文形成三项有界贡献：（1）形式化给出 \proto{}/\fragproto{} 帧布局、解析式提交波形成本模型以及确定性包内标记偏移，并通过有线馈线实采完成验证；（2）通过 36 轮交替区组无线实验（\num{3600} 项计划任务），在相等单次尝试波形时长下隔离共包带来的突发次数缩减（\modejc{} 对 \modesc{}，同用 \ble{} ch3 / \stdfour{} ch12），并定量揭示相对 ch39 \modelc{} 由载荷与预算策略驱动的单位成本排序翻转及分片损益；（3）评估固定次信道时间可用性与单侧断馈恢复行为（第~\ref{sec:design}--\ref{sec:eval} 节、图~\ref{fig:layout}--\ref{fig:p4_tradeoff}、表~\ref{tab:budget}--\ref{tab:p4_main}）。
